\subsection{Linear Prediction and Generalized Linear Models}

Linear prediction reduces inputs to scores of the form ``intercept plus weighted features.'' Regression reads scores as predictions; classification compares or transforms them.  Its appeal is geometric: projections, hyperplanes, and feature maps.

\begin{definition} \label{def:sl-linear-predictor} \hlt{Linear Predictor}\\
For $x\in\R^p$, $w\in\R^p$, and $b\in\R$, the \hlt{linear predictor} is $f_{w,b}(x)\defeq w^\top x+b$.\\
Regression uses $f_{w,b}$ directly; classification transforms it by $g$. Score level sets are parallel hyperplanes.
\end{definition}

\begin{definition} \label{def:sl-least-squares-linear-regression} \hlt{Least-Squares Linear Regression}\\
For $X\in\R^{n\times d}$ and $y\in\R^n$, \hlt{least-squares linear regression} chooses $\widehat{\beta}\in\operatorname*{argmin}\limits_{\beta\in\R^d}\norm{y-X\beta}_2^2$ and $\widehat{y}\defeq X\widehat{\beta}$.\\
If $X^\top X$ is nonsingular, then $\widehat{\beta}=(X^\top X)^{-1}X^\top y$. Columns of $X$ may be raw variables, transformations, indicators, interactions, or basis functions; linearity is in $\beta$.
\end{definition}

Geometrically, least squares moves $y$ to the closest point in $\operatorname{col}(X)$.

\begin{method} \label{meth:sl-linear-regression-projection-and-rank} \hlt{Projection View of Least Squares}
\begin{enumerate}[label=(\roman*)]
\setlength{\itemsep}{0pt}
\item \hlt{Full rank}: If $\operatorname{rank}(X)=d$, define $H\defeq X(X^\top X)^{-1}X^\top$ and $\widehat{y}\defeq Hy$.
\item \hlt{Projection}: $H$ is the orthogonal projection onto $\operatorname{col}(X)$; hence $y-\widehat{y}\perp\operatorname{col}(X)$.
\item \hlt{Rank deficiency}: If $\operatorname{rank}(X)<d$, $\widehat{\beta}$ may be nonunique, but $\widehat{y}=\operatorname{proj}_{\operatorname{col}(X)}y$ is unique.
\end{enumerate}
Rank deficiency affects coefficient representation, not least-squares fitted values.
\end{method}

Successive orthogonalization fits in an orthogonal basis for the same column space.

\begin{breakablealgorithm}
\caption{Least Squares by Successive Orthogonalization}
\label{alg:sl-successive-orthogonalization-least-squares}
\begin{algorithmic}
\Require Full-column-rank design columns $x_0,\ldots,x_p$ and response vector $y$
\Ensure Least-squares coefficients $\widehat{\beta}$
\For{$j=0,\ldots,p$}
\State $z_j\gets x_j$
\For{$\ell=0,\ldots,j-1$}
\State $r_{\ell j}\gets q_\ell^Tx_j$
\State $z_j\gets z_j-r_{\ell j}q_\ell$
\EndFor
\State $r_{jj}\gets \norm{z_j}_2$
\State $q_j\gets z_j/r_{jj}$
\EndFor
\State Form $Q=[q_0,\ldots,q_p]$ and the upper-triangular matrix $R=(r_{\ell j})$
\State Solve $R\widehat{\beta}=Q^Ty$ by back-substitution
\end{algorithmic}
\end{breakablealgorithm}

\begin{remark} \label{rem:sl-linear-inference-after-model-choice} \hlt{Linear-Model Inference After Selection}\\
Under fixed-design $y=X\beta+\varepsilon$, $\varepsilon\sim N(0,\sigma^2I_n)$, and $\operatorname{rank}(X)=d$, $\widehat{\beta}\sim N(\beta,\sigma^2(X^\top X)^{-1})$ and $\widehat{\sigma}^2\defeq\frac{\norm{y-X\widehat{\beta}}_2^2}{n-d}$.\\
These give standard errors, $t$ or asymptotic $z$ statistics, and nested-model $F$ tests when $X$ is fixed in advance.\\
After data-driven selection of $X$ or the model, these standard errors ignore selection.
\end{remark}

\begin{method} \label{meth:sl-multiple-output-least-squares} \hlt{Multiple-Output Least Squares}\\
For $Y\in\R^{n\times K}$ and $B\in\R^{d\times K}$, set $Y=XB+E$.
\begin{enumerate}[label=(\roman*)]
\setlength{\itemsep}{0pt}
\item \hlt{Least squares}: If $\operatorname{rank}(X)=d$, then $\widehat{B}=(X^\top X)^{-1}X^\top Y$.
\item \hlt{Decoupling}: Writing $Y=(y_1,\ldots,y_K)$ gives $\widehat{B}=(\widehat{\beta}_1,\ldots,\widehat{\beta}_K)$, where each $\widehat{\beta}_k$ is the OLS fit of $y_k$ on $X$.
\end{enumerate}
Shared responses can be pooled by low-rank $B$ or common-direction shrinkage.
\end{method}

\begin{method} \label{meth:sl-canonical-correlation-and-reduced-rank-regression} \hlt{Canonical Correlation and Reduced-Rank Regression}
\begin{enumerate}[label=(\roman*)]
\setlength{\itemsep}{0pt}
\item \hlt{CCA}: Choose $(v_j,u_j)$ with large $\operatorname{corr}(Xv_j,Yu_j)$, orthogonal across $j$; at most $\min(d,K)$ pairs.
\item \hlt{Rank constraint}: Fit $Y\approx XB$ with $\operatorname{rank}(B)\leq m$, restricting fitted responses to $m$ dimensions.
\item \hlt{Objective}: $\widehat{B}^{\mathrm{rr}}_m\in\operatorname*{argmin}\limits_{\operatorname{rank}(B)\leq m}\sum\limits_i(y_i-B^\top x_i)^\top\widehat{\Sigma}^{-1}(y_i-B^\top x_i)$.
\item \hlt{Role of $m$}: $m=K$ gives separate OLS; smaller $m$ pools shared response structure.
\end{enumerate}
Reduced rank assumes the predictable response variation lies in a low-dimensional subspace.
\end{method}

\begin{definition} \label{def:sl-discriminant-functions-linear-boundaries} \hlt{Discriminant Functions and Linear Boundaries}\\
For $\mathcal{G}=\{1,\ldots,K\}$ and scores $\delta_k:\mathcal{X}\to\R$, set $\widehat{G}(x)\in\operatorname*{argmax}_{k\in\mathcal{G}}\delta_k(x)$.  Only score differences matter.
\begin{enumerate}[label=(\roman*)]
\setlength{\itemsep}{0pt}
\item \hlt{Pairwise boundary}: The $k$ versus $\ell$ boundary lies in $\{x:\delta_k(x)=\delta_\ell(x)\}$.
\item \hlt{Linear boundary}: If $\delta_k-\delta_\ell$ is affine in $x$ for all $k,\ell$, then all pairwise decision boundaries are hyperplanes.
\item \hlt{Binary log-odds}: If $\log\frac{\mathbb{P}(G=1\mid X=x)}{\mathbb{P}(G=2\mid X=x)}=\beta_0+\beta^\top x$, then the boundary is $\{x:\beta_0+\beta^\top x=0\}$.
\end{enumerate}
\end{definition}

\begin{method} \label{meth:sl-indicator-matrix-regression-classifier} \hlt{Indicator-Matrix Regression Classifier}\\
For classes $g_i\in\{1,\ldots,K\}$, set $t_k=e_k\in\R^K$, $Y_i=t_{g_i}$, and $\widetilde{x}_i^\top=(1,x_i^\top)$, coding labels as Euclidean targets.
\begin{enumerate}[label=(\roman*)]
\setlength{\itemsep}{0pt}
\item \hlt{Least-squares fit}: Choose $\widehat{B}\in\operatorname*{argmin}_B\sum\limits_{i=1}^{n}\norm{Y_i-B^\top\widetilde{x}_i}_2^2$; if $\widetilde{X}^\top\widetilde{X}$ is nonsingular, $\widehat{B}=(\widetilde{X}^\top\widetilde{X})^{-1}\widetilde{X}^\top Y$.
\item \hlt{Scores}: For a new $x$, set $\widetilde{x}=(1,x^\top)^\top$ and $\widehat{f}(x)\defeq\widehat{B}^\top\widetilde{x}$.
\item \hlt{Classification}: Predict $\widehat{G}(x)\in\operatorname*{argmax}_k\widehat{f}_k(x)$, equivalently nearest $t_k$.
\end{enumerate}
With an intercept, $\sum_k\widehat{f}_k(x)=1$, but $\widehat{f}_k(x)$ may leave $[0,1]$ and, for $K\geq3$, \hlt{mask} middle classes.
\end{method}

\begin{definition} \label{def:sl-generalized-linear-model} \hlt{Generalized Linear Model}\\
A \hlt{generalized linear model} gives $Y\mid X=x$ an exponential-family law and makes a linked mean linear:\\
$g(\mu(x))=w^\top x+b$. Logistic regression uses the Bernoulli law and logit link $g(\mu)=\log\frac{\mu}{1-\mu}$.
\end{definition}

Thus a GLM separates score, mean scale, and likelihood.

\begin{definition} \label{def:sl-multiclass-logistic-regression} \hlt{Multiclass Logistic Regression}\\
With class $K$ as baseline, set $\eta_k(x)\defeq\beta_{k0}+\beta_k^\top x=\log\frac{\mathbb{P}(G=k\mid X=x)}{\mathbb{P}(G=K\mid X=x)}$ for $k<K$.
\begin{enumerate}[label=(\roman*)]
\setlength{\itemsep}{0pt}
\item \hlt{Softmax probabilities}: For $k<K$, $p_k(x;\theta)\defeq\frac{\exp(\eta_k(x))}{1+\sum\limits_{\ell=1}^{K-1}\exp(\eta_\ell(x))}$, and $p_K(x;\theta)\defeq\frac{1}{1+\sum\limits_{\ell=1}^{K-1}\exp(\eta_\ell(x))}$.
\item \hlt{Likelihood}: Fit $\theta$ by maximising $\ell(\theta)\defeq\sum\limits_{i=1}^{n}\log p_{g_i}(x_i;\theta)$.
\end{enumerate}
The softmax gives $p(x;\theta)\in\Delta^{K-1}$; the reference class changes only the parametrisation.
\end{definition}

\begin{method} \label{meth:sl-logistic-irls} \hlt{Binary Logistic Regression by IRLS}\\
For $y_i\in\{0,1\}$ and intercept-augmented rows $x_i^\top$, set $p_i(\beta)\defeq(1+\exp(-x_i^\top\beta))^{-1}$.
\begin{enumerate}[label=(\roman*)]
\setlength{\itemsep}{0pt}
\item \hlt{Derivatives}: With $p=(p_i)$ and $W\defeq\operatorname{diag}(p_i(1-p_i))$, $\nabla\ell(\beta)=X^\top(y-p)$ and $\nabla^2\ell(\beta)=-X^\top W X$.
\item \hlt{Adjusted response}: At $\beta^{\mathrm{old}}$, set $z\defeq X\beta^{\mathrm{old}}+W^{-1}(y-p)$.
\item \hlt{IRLS update}: Set $\beta^{\mathrm{new}}\defeq(X^\top W X)^{-1}X^\top Wz$.
\end{enumerate}
IRLS is Newton's method as weighted least squares.\\
The log-likelihood is concave when the MLE is finite; step-halving controls Newton overshoot.
\end{method}

\begin{method} \label{meth:sl-l1-logistic-regression} \hlt{$L_1$-Regularized Logistic Regression}\\
With unpenalized intercept, maximize $\ell_\lambda(\beta_0,\beta)\defeq\sum\limits_{i=1}^{n}\{y_i\eta_i-\log(1+\exp(\eta_i))\}-\lambda\sum\limits_{j=1}^{p}\abs{\beta_j}$, where $\eta_i\defeq\beta_0+x_i^\top\beta$.
\begin{enumerate}[label=(\roman*)]
\setlength{\itemsep}{0pt}
\item \hlt{Probabilities}: Set $p_i\defeq(1+\exp(-\eta_i))^{-1}$ and residual $r\defeq y-p$.
\item \hlt{KKT conditions}: If $\widehat{\beta}_j\neq0$, then $x_j^\top r=\lambda\sign(\widehat{\beta}_j)$; if $\widehat{\beta}_j=0$, then $\abs{x_j^\top r}\leq\lambda$.
\end{enumerate}
The KKT conditions are thresholded feature--residual correlations.
\end{method}

\begin{method} \label{meth:sl-gaussian-discriminant-analysis} \hlt{Gaussian Discriminant Analysis}\\
If $X\mid G=k$ has density $f_k$ and prior $\pi_k$, then $\mathbb{P}(G=k\mid X=x)=\frac{f_k(x)\pi_k}{\sum\limits_{\ell=1}^{K}f_\ell(x)\pi_\ell}$.
\begin{enumerate}[label=(\roman*)]
\setlength{\itemsep}{0pt}
\item \hlt{Score}: If $X\mid G=k\sim N(\mu_k,\Sigma_k)$, compare $\delta_k(x)\defeq-\frac{1}{2}\log|\Sigma_k|-\frac{1}{2}(x-\mu_k)^\top\Sigma_k^{-1}(x-\mu_k)+\log\pi_k$.
\item \hlt{LDA}: If $\Sigma_k=\Sigma$ for all $k$, equivalent linear scores are $\delta_k^{\mathrm{LDA}}(x)\defeq x^\top\Sigma^{-1}\mu_k-\frac{1}{2}\mu_k^\top\Sigma^{-1}\mu_k+\log\pi_k$.
\item \hlt{QDA}: If $\Sigma_k$ varies with $k$, the quadratic term remains and the rule is quadratic discriminant analysis.
\end{enumerate}
Estimate $\pi_k,\mu_k,\Sigma_k$ from training data; pooled covariance gives LDA and classwise covariances give QDA.\\
Shared covariance makes the boundaries linear.
\end{method}

\begin{method} \label{meth:sl-regularized-discriminant-analysis} \hlt{Regularized Discriminant Analysis}\\
For $0\leq\alpha\leq1$, set $\widehat{\Sigma}_k(\alpha)\defeq\alpha\widehat{\Sigma}_k+(1-\alpha)\widehat{\Sigma}$.
\begin{enumerate}[label=(\roman*)]
\setlength{\itemsep}{0pt}
\item \hlt{Endpoints}: $\alpha=0$ gives LDA and $\alpha=1$ gives QDA.
\item \hlt{Scalar shrinkage}: For $0\leq\gamma\leq1$, replace $\widehat{\Sigma}$ by $\widehat{\Sigma}(\gamma)\defeq\gamma\widehat{\Sigma}+(1-\gamma)\widehat{\sigma}^{2}I$.
\end{enumerate}
Shrinkage reduces variance when covariance estimation is unstable.
\end{method}

\begin{definition} \label{def:sl-naive-bayes-classifier} \hlt{Naive Bayes Classifier}\\
For $x=(x_1,\ldots,x_p)$, assume $f_j(x)=\prod\limits_{k=1}^{p}f_{jk}(x_k)$ for each class $j$.
\begin{enumerate}[label=(\roman*)]
\setlength{\itemsep}{0pt}
\item \hlt{Posterior}: With priors $\pi_j$, $p_j(x)\defeq\mathbb{P}(G=j\mid X=x)\propto\pi_j\prod\limits_{k=1}^{p}f_{jk}(x_k)$.
\item \hlt{Additive log odds}: Relative to class $J$, $\log\frac{p_\ell(x)}{p_J(x)}=\log\frac{\pi_\ell}{\pi_J}+\sum\limits_{k=1}^{p}\log\frac{f_{\ell k}(x_k)}{f_{Jk}(x_k)}$.
\end{enumerate}
On the log scale, features contribute additively.\\
Even misspecified independence can reduce variance enough to classify well near boundaries.
\end{definition}

\begin{remark} \label{rem:sl-least-squares-versus-lda} \hlt{Least Squares Versus LDA}\\
For $K=2$, least-squares regression on any distinct label coding gives $\widehat{\beta}\propto\widehat{\Sigma}^{-1}(\widehat{\mu}_2-\widehat{\mu}_1)$, the LDA direction.\\
The cutpoint may differ unless classes are balanced or tuned separately.\\
For $K>2$, LDA uses pooled-Mahalanobis centroid distances and avoids masking.
\end{remark}

\begin{remark} \label{rem:sl-logistic-versus-lda} \hlt{Logistic Regression Versus LDA}\\
Both can yield linear log-odds. Logistic regression models class probabilities directly and leaves $P_X$ free; LDA fits a joint Gaussian mixture for $(X,G)$ and uses Bayes' rule.
\end{remark}

\begin{remark} \label{rem:sl-lda-whitening-and-reduced-rank-view} \hlt{Whitening and Reduced-Rank LDA}\\
With common covariance $W$, LDA compares $W^{-1/2}x$ to whitened centroids, adjusted by $\log\pi_k$.\\
The whitened centroids span dimension at most $K-1$, so LDA reduces to this centroid subspace.\\
Fisher directions maximize $\frac{a^\top Ba}{a^\top Wa}$, equivalently solve $Ba=\lambda Wa$, where $B$ is between-class scatter.
\end{remark}

\begin{remark} \label{rem:sl-feature-maps-linear-models} \hlt{Feature Maps}\\
A feature map $\phi:\mathcal{X}\to\mathcal{H}$ gives linear prediction in feature space, $f(x)=w^\top\phi(x)+b$, which is generally nonlinear in $x$. Kernel methods choose $\phi$ implicitly; neural networks learn $\phi$ from data.
\end{remark}
